% 本文件由 LaTeX 排版 Agent 根据有效 translation.json 自动生成。
% author=Athanasius; work_id=2809; title=论尼西亞大公会议法令

\chapter{论\Place{尼西亞}大公会议法令}

% structure: p0001 source=h1 target=omitted reason=page-title-already-rendered-by-parent

% structure: p0002 source=h2 target=section reason=relative-heading-rank; base=h2

\section{第一章 导言。\Emph{亚略派反对\Place{尼西亞}大公会议的控诉；他们的反复无常；他们如同犹太人；他们使用武力而非理性。}}

1. 你做得很好，向我通报了你与亚略派拥护者的讨论，其中有些是\Person{欧瑟伯}的朋友，还有许多持守教会教义的弟兄。我因你对基督的爱而赞许你的警醒，这警醒卓越地揭露了他们异端的无宗教性；同时我惊异于亚略派的无耻，在过去他们的论据已被揭露为错误和虚浮，甚至众人普遍认为他们极其乖谬之后，他们仍然像犹太人一样抱怨，\Quote{为什么\Place{尼西亞}的教父们使用圣经中没有的术语，即\InnerQuote{\OriginalTerm{出于本质}{\GreekText{ἐκ} \GreekText{τῆς} \GreekText{οὐσίας}}}和\InnerQuote{\OriginalTerm{同一本质}{\GreekText{ὁμοούσιος}}}？}。而你，作为一个有学问的人，不顾他们的借口，指证他们说话毫无意义；而他们设词遮掩，正适合他们邪恶的性情。因为他们情绪多变不定，就像变色龙的颜色；当被揭露时，他们显得慌乱，被提问时他们犹豫，然后他们失去羞耻，转而逃避。当他们在这方面被识破，他们不肯休息，直到编造出新的不存在的事，并且，按照圣经，\BibleQuote{图谋虚妄}；这一切都是为了坚持他们的无宗教性。\par

如今这种努力无非是他们缺乏理性的明显标志，而且如我所说，是效法犹太人的恶毒。因为犹太人也是这样，当被真理指责，无法对抗时，就使用借口，例如，\BibleQuote{你行什么神迹，叫我们看见就信你？你到底作什么事呢？}\BibleRef{若 6:30}。尽管有那么多神迹赐下，他们自己却说，\Quote{我们怎么办呢？因为这人行许多奇迹。}事实上，死人复活，瘸子行走，瞎子重见光明，癞病人得了洁净，水变成酒，五个饼使五千人吃饱，众人都惊奇并敬拜主，承认在祂身上应验了预言，祂是天主，天主子；只有法利塞人，虽然神迹比太阳还明亮，却仍然像无知的人抱怨，\Quote{你既是人，为什么把自己当作天主呢？}愚蠢，且真是心灵盲目！他们反应该说，\Quote{你既是天主，为什么成为人？}因为祂的作为证明祂是天主，好叫他们既崇拜圣父的美善，又赞叹圣子为我们缘故所行的\OriginalTerm{经世}{\GreekText{οἰκονομία}}。然而他们不说这话；不，他们也不喜欢观看祂所做的事；或者他们确实看见了，因为他们无法避免，但他们又改变了抱怨的理由，\Quote{你为什么在安息日治好瘫子，为什么叫胎生的瞎子看见？}但这也不过是借口，纯粹的抱怨而已；因为在别的日子主也医治\BibleQuote{各种病症，各种疾患}\BibleRef{玛 4:23}，但他们照常抱怨，并称祂为别西卜，宁愿担上不信神的嫌疑，也不肯弃绝自己的邪恶。尽管救主在这样不同的时间和各种方式中彰显了祂的\OriginalTerm{天主性}{\GreekText{θεότης}}，并向众人宣讲了圣父，然而，他们如同用脚踢刺，以愚昧的语言反对，他们这样做，正如神圣箴言所说，以便寻找借口，使自己与真理隔绝。\par

2. 正如那时代的犹太人因为如此邪恶行事并否认主，就公义地被剥夺了他们的法律和向他们祖先所作的应许，同样，如今犹太化的亚略派，依我判断，正处于\Person{盖法}和当代法利塞人的境地。因为，他们发觉自己的异端完全不合理，便捏造借口，\Quote{为什么这项定了，那项却不定义呢？}然而，他们现在如此行事不要惊异；因为不久他们就会转为暴力，接着就会威胁\Quote{差役和官长}。确实，正如我们所见，他们的异端思想就在这些事上得到支持；因为他们否认天主的\OriginalTerm{圣言}{\GreekText{Λόγος}}，完全没有理性，这也是公义的。既意识到这一点，我本不想回答他们的询问；但是，既然你的友善请求了解大公会议的经过，我便立刻毫不迟延地叙述了当时发生的事，用简短的话指出，亚略主义完全缺乏宗教精神，他们唯一的事就是制造借口。\par

% structure: p0006 source=h2 target=section reason=relative-heading-rank; base=h2

\section{第二章 亚略派面对\Place{尼西亞}大公会议的行径。\Emph{既无知又不虔敬，竟企图表决推翻大公会议：\Place{尼西亞}的会议进程：\Person{欧瑟伯}派当时签署了他们如今抱怨的文件：论真教师的同心合意与传统的过程：亚略派的变化。}}

而你，亲爱的，请考虑是否如此。如果，魔鬼在他们的心中播下了这种乖谬，他们便对自己邪恶的构想充满信心，那么让他们针对已提出的异端证据为自己辩护，然后，如果可能的话，再来挑剔针对他们而制定的定义。因为没有人因谋杀或奸淫被判有罪之后，还能在审判后指责法官的判决，说为什么他这样说，而不那样说。因为这不能为罪犯开脱，反而以其任性大胆而加重了他的罪行。同样，让这些人要么证明他们的思想是虔诚的（因为他们当时被告且被定罪，他们的抱怨是后来的，受指控的人理应只为自己辩护），要么如果他们良心不洁，意识到自己不虔敬，就不该抱怨他们不理解的事，否则他们将给自己带来双重指责：不虔敬与无知。不如让他们以温顺的精神探究此事，学习他们至今未知的事，用真理之泉和宗教教义洗净他们不虔敬的耳朵。\par

3. 现在，在尼西亚大公会议上发生在\Person{欧瑟比}及其同党身上的事如下：——当他们坚持自己的不虔，并试图与天主作战时，他们所采用的言辞充满了不虔；但聚集在一起的三百多位主教，温和而慈爱地要求他们在宗教立场上作出解释并为自己辩护。然而，他们几乎刚一开口，就被定罪，而且彼此意见分歧；他们意识到自己的异端陷入了困境，便闭口无言，用沉默承认了临到他们异说上的羞耻。于是，主教们否定了他们编造的言辞，宣布了反对他们的纯正且符合教会传统的信仰；而当所有人都签署这份信仰时，\Person{欧瑟比}及其同党也签署了同样的言辞，也就是他们现在所抱怨的那些话，我指的是 \Quote{出于\OriginalTerm{性体}{essence}}、\Quote{\OriginalTerm{同性体}{one in essence}}，以及 \Quote{天主圣子既非受造物，亦非被造者，不属受造之列，而是出自圣父\OriginalTerm{性体}{substance}的圣言。} 而更奇特的是，前一天还加以否认、后来却签署了的巴勒斯坦\Place{凯撒勒亚}的\Person{欧瑟比}，写信给他所属的教会说，这便是教会的信仰和教父们的传承；并且公开承认他们之前犯了错误，曾鲁莽地对抗真理。尽管他当时羞于采用这些说法，并以自己的方式向教会辩解，然而在他的书信中，由于他没有否认 \Quote{同性体} 和 \Quote{出于性体}，他无疑暗示了所有这一切。这样一来，他就陷入了困境；因为就在为自己辩解的同时，他进而攻击亚略派，指责他们声称 \Quote{圣子在祂受生之前并不存在}，并由此否认祂在降生成人之前的存在。而\Person{阿卡西乌}也知道此事，尽管他同样可能由于时局而心存畏惧，假装不知并否认事实。因此，我在末尾附上了\Person{欧瑟比}的信，好叫你们从中得知，基督的仇敌对他们自己的导师表现出了何等的不敬，尤其是\Person{阿卡西乌}本人。\par

4. 难道他们不是在思想上反对如此伟大而大公的会议，就是犯了罪吗？当他们竟敢对抗那反对亚略主义的善美定义，即甚至连当初教给他们不虔的人也已承认的定义，难道不是犯了过犯吗？即使签署之后，\Person{欧瑟比}及其同党确实再次改变，像狗转过来吃自己所吐的不虔，难道眼前的这些反对者不更应受憎恶吗？因为他们如此将自己的灵魂自由牺牲于他人，情愿以这些人为他们异端的主子，而正如\Person{雅各伯}所说的，他们是 \BibleQuote{三心两意的人，在他一切的行径上，易变无定。}\BibleRef{雅 1:8} 他们没有一个固定的意见，而是反复无常，一会儿推荐某些说法，一会儿又贬低它们，反过来推荐刚才还在指责的东西。然而，正如\WorkTitle{牧人书}所言，这是 \Quote{魔鬼之子}，是商贩而非导师的特征。因为，我们的教父们所传授的，才是真正的教义；而这才是导师的真正标志：彼此承认同一件事，既不背离自己，也不背离自己的教父；但凡没有这种品格的人，不可称为真导师，而是恶的。因此，希腊人由于没有见证相同的教义，反而彼此争吵，便没有教导的真理；但神圣而真实的真理宣报者们却和谐一致，没有分歧。尽管他们生活在不同的时代，却全都同道而行，是唯一天主的先知，和谐一致地宣讲同一圣言。\par

5. 因此，\Person{梅瑟}所教导的，\Person{亚巴辣罕}遵守了；\Person{亚巴辣罕}所遵守的，\Person{诺厄}和\Person{哈诺客}也认明了，他们能分别洁净与不洁，并成为蒙天主悦纳的人。因为\Person{亚伯尔}也这样作了见证，他知道他从\Person{亚当}所学到的，而\Person{亚当}是从那位主领受的；主在末世来临，为除灭罪过时曾说：\BibleQuote{我赐给你们的不是一条新诫命，而是一条旧诫命，就是你们从起初所听过的。}\BibleRef{若一 2:7}因此，\OriginalTerm{真福宗徒}{blessed Apostle}\Person{保禄}，既从祂领受了这道理，在描述教会职务时，也禁止执事——更不用说主教——一口两舌\BibleRef{弟前 3:8}；并且他在谴责\Place{迦拉达}人时，作了明确的声明：\BibleQuote{无论是谁，即使是我们，或是从天上降下的一位天使，若给你们宣讲的福音，与我们给你们所宣讲的福音不同，当受绝罚。正如我们以前说过的，如今我再说一次：谁若给你们宣讲福音与我们给你们所宣讲的不同，当受绝罚。}\BibleRef{迦 1:8}既然宗徒如此说，那么这些人或是应当绝罚\Person{欧瑟比乌}及其同党，至少因为他们反复无常，宣认与他们所签署的相反的道理；或是，如果他们承认自己的签署是正当的，就不应对如此伟大的大公会议口出怨言。但如果他们两者都不做，他们就太明显是随从各种异风与浪潮的玩物，受别人的意见所左右，而不是自己的；既是如此，他们如今所主张的，也像从前一样不值得尊重。最好让他们停止非议他们所不明白的；免得因不辨是非，竟称恶为善，称善为恶，以苦为甜，以甜为苦。无疑，他们渴望那些已被判为错误并遭弃绝的教义得以占上风，并且竭力歪曲已正确定断的道理。我们实不必再作进一步的解释或答复他们的借口，他们也不应再顽抗，而是该默许他们异端首领所签署的；因为，\Person{欧瑟比乌}及其同党后来的转变虽可疑而不道德，但他们当时至少还有一点为自己辩护的机会时所作的签署，却确证了他们教义的不虔。因为他们若非曾谴责这异端，先也不会签署；若非陷入困境与羞耻，也不会谴责它；因此，再度转变正证明他们对不虔的好争之心。因此，正如我所说过的，这些人也该安分守己；但是，既然他们由于极度的放肆，或许指望能比别人更好地拥护这种魔鬼般的不虔，所以，尽管在前一封写给你们的信中，我已详细驳斥过他们，然而，来吧，让我们现在也逐一检查他们个别的说法，如同对他们的前辈一样；因为如今并不亚于当初，他们的异端必被证明毫不可靠，而是出于恶神。\par

% structure: p0011 source=h2 target=section reason=relative-heading-rank; base=h2

\section{第三章 \Emph{\OriginalTerm{子}{Son}这个词的两种意义：1. \OriginalTerm{嗣子}{adoptive}; 2. \OriginalTerm{本体的}{essential}; 亚略派试图在两者之间寻找第三种意义；例如，主只是由天主直接受造（\Person{亚斯特留}的观点），或唯独主分享了圣父。第二种且真实的意义；天主怎样创造，也怎样真实地生；虽然祂的创造与\OriginalTerm{受生}{generation}不同于人的；祂的受生不依赖时间；受生意味着天主内一种内在的，因而也是永恒的行动；对 \BibleBook{}{箴言}8:22 的解释。}}

6. 于是他们说，这正是别人以前所持守、并胆敢主张的：\Quote{不是常存之父，不是常存之子；因为子受生之前并不存在，而是像其他受造物一样，从虚无中产生；因此，天主并不常是子之父；而是当子产生并被造时，天主才被称为祂的父亲。因为圣言是一个受造物和工程，在本质上与圣父疏远且不相似；子既非按本性是圣父的真实圣言，也不是祂的独一而真实的智慧；而是一个受造物和工程之一，不当被称为圣言和智慧；因为祂是由那寓于天主内的圣言所造的，如同万有一样。所以，子不是真天主。}\par

现在，为了让他们明白自己所说的是什么，最好先问他们这一点：事实上，\OriginalTerm{儿子}{son}是什么，这个名称表示什么。的确，天主的圣经使我们知晓这词的两种意义：——一种是\Person{梅瑟}在法律书中向我们提出的，\BibleQuote{当你们听从上主你们天主的声音，遵守我今天吩咐你们的一切诫命，做上主你们天主眼中视为正直的事时，你们就是上主你们天主的子女}；同样，在福音中，\Person{若望}说：\BibleQuote{凡接受祂的，祂赐给他们权能，好成为天主的子女。}\BibleRef{若 1:12}——另一种意义是：\Person{依撒格}是\Person{亚巴辣罕}的儿子，\Person{雅各伯}是\Person{依撒格}的儿子，诸位圣祖是\Person{雅各伯}的儿子。那么，他们讲论天主之子，讲出前面那些无稽之谈时，是在哪一种意义上理解的呢？因为我确信，他们将与\Person{欧瑟比乌}及其同党一样，陷入同样的不虔。\par

如果是在第一种意义上，就是那些由于道德进步而藉恩宠获得这名称，并领受权能成为天主子女的人（因为这是他们前辈所说的），那么祂就似乎与我们毫无区别了；不，祂也不是独生子，因为祂按功德像其他人一样获得了儿子的名分。因为，即使允许他们所说，由于祂的资质被预先知道，因此祂从起初，从祂存在的第一刻，就获得了这恩宠，这名称和这名称的光荣，祂与那些在行为后才获得这名称的人仍然没有区别，只要这也是祂像其他人一样取得儿子身份的基础。因为\Person{亚当}虽然从起初就领受了恩宠，一受造便立刻被安置在乐园里，但他与\Person{哈诺客}没有分别——\Person{哈诺客}在出生后一段时间才因中悦天主而被提升到乐园——，也与那位行为后被提到乐园的宗徒没有分别；甚至与那曾是强盗的人也没有分别，他因宣认，获得应许，立刻进入乐园。\par

7. 当他们受到这样的追问时，或许会给出一个已多次给他们带来麻烦的回答； \Quote{我们认为，圣子拥有超越其他一切的特权，因此被称为独生子，因为唯独祂是由天主独自所生，而其他万物都是由天主藉着圣子所造。} 如今我纳闷，是谁向你们提出了如此荒谬而新奇的观念，以为圣父独自用祂自己的\OriginalTerm{手}{Hand}造了圣子，而其他万物都是由圣子如同一个帮手所造。倘若由于辛劳的缘故，天主仅满足于创造圣子，而非立即创造万物，那么这是一个不虔敬的思想，尤其是对那些知道\Person{依撒意亚}先知的话的人，\BibleQuote{永恒的天主，上主，大地的创造者，不饥不倦；祂的智慧深不可测。}\BibleRef{依 40:28} 相反，是祂赐力量给饥饿的人，并藉着祂的\OriginalTerm{圣言}{Word}使劳苦者得享安息。再者，认为祂鄙弃亲自造那些在圣子之后的受造物，视其为卑微的任务，这也是不虔敬的；因为那位天主毫无骄傲，祂与\Person{雅各伯}一同下到\Place{埃及}，为了\Person{亚巴郎}的缘故，因\Person{撒辣}而责罚\Person{阿彼默肋客}，与身为人的\Person{梅瑟}面对面说话，降临在\Place{西乃山}上，并以祂隐秘的恩宠为祂的百姓与\Person{阿玛肋克}作战。然而，你们的这个说法也是错误的，因为\BibleQuote{祂创造了我们，不是我们自己。} 是祂藉着自己的\OriginalTerm{圣言}{Word}创造了万物，无论大小；我们不可将创造分割，说这是圣父的，那是圣子的；它们都属于同一天主，祂使用自己的\OriginalTerm{圣言}{Word}如同\OriginalTerm{手}{Hand}，并在祂内作成万事。这位天主亲自向我们显示，当祂说：\BibleQuote{这一切都是我的手所造的}\BibleRef{依 66:2}；而\Person{保禄}宗徒按他所领受的教导我们：\BibleQuote{只有一个天主，万物都出于祂；也只有一个主耶稣基督，万物都藉着祂而存在。}\BibleRef{格前 8:6} 因此，祂自始至终都是如此：祂吩咐太阳，太阳便升起；命令云彩，雨便降在一处；雨不降的地方，就变成干旱。祂命大地出产果实，并在母胎中形成了\Person{耶肋米亚}\BibleRef{耶 1:5}。如果祂现今行这一切，那么在起初，祂也必定不会鄙弃亲自藉着\OriginalTerm{圣言}{Word}创造万物；因为这些不过是整体中的部分。\par

8. 然而，我们暂且假设其他受造物无法承受由\OriginalTerm{无始者}{Unoriginate}的直接之\OriginalTerm{手}{Hand}所造就，因而唯有圣子由圣父独自所生，其他万物则由圣子如同帮手与助理所造——因为这就是献祭者\Person{阿斯特里乌斯}所写的，而\Person{亚略}加以抄录并传给了自己的朋友们，从那时起他们便使用这种言辞，它如同压伤的芦苇，这些糊涂人却不知它是何等地脆弱。因为，如果\OriginalTerm{有始的}{originate}之物不可能承受天主的\OriginalTerm{手}{Hand}，而你们认为圣子也是其中之一，那么祂又怎能堪当由天主独自造就呢？再者，如果为了有始之物得以存在，必须有一位\OriginalTerm{中保}{Mediator}，而你们认为圣子也是有始的，那么，在祂之前必定另有某位中介，为祂的受造而存在；那位中保本身又是受造物，依此类推，他自己也需要另一位中保来构成自己。即使我们再设想另外一位，就必须先设想他的中保，这样我们便永无止境。于是，由于中保的需求不断，受造界就永远无法建立，因为按你们所说，任何有始之物都不能承受\OriginalTerm{无始的}{Unoriginate}的直接之\OriginalTerm{手}{Hand}。如果你们意识到这种荒谬，便开始说：圣子虽是受造物，却被造得能够由\OriginalTerm{无始者}{Unoriginate}所造；那么，其它事物也必定同样能够由\OriginalTerm{无始者}{Unoriginate}直接所造，因为按你们的判断，圣子也只不过是万有中的一个受造物罢了。因此，按照你们不虔敬而又虚妄的想象，\OriginalTerm{圣言}{Word}的受造就成了多余；因为天主足以直接形成万物，而所有有始之物也都能承受祂的直接之\OriginalTerm{手}{Hand}。\par

这些不虔敬的人，如此疯狂而心智短浅，让我们看看这个特殊的诡辩是否比其他的更不合理。\Person{亚当}乃是天主独自藉着\OriginalTerm{圣言}{Word}所造的唯一之人；然而，没有人会说亚当比其他人更优越，或与他的后代不同，尽管他是由天主独自创造和塑造的，而我们都是亚当的后裔，依种族血脉繁衍；因为他像其他人一样由泥土造成，起初并不存在，后来才得以存在。\par

9. 但即便我们承认\OriginalTerm{原祖}{Protoplast}因曾被视为被天主亲手所造而享有某种特权，这特权也只是尊荣上的，而非本性上的。因为他是由土而出的，如同其他人一样；而那时形成亚当的手，也同样在现在和永远形成并赋予他后来的人完整的存有。天主亲自对\Person{耶肋米亚}说明了这一点，正如我之前说过的；\BibleQuote{我还没有在母腹内形成你以前，我已认识了你}\BibleRef{耶 1:5}；同样，祂论及一切时说：\BibleQuote{这一切都是我手所造的}\BibleRef{依 66:2}；又藉\Person{依撒意亚}说：\BibleQuote{你的救赎主，由母腹中形成你的上主这样说：我是创造一切的上主，独自展开了诸天，铺开了大地。} 于是\Person{达味}明白这一点，便在\BibleBook{}{圣咏}中说：\BibleQuote{你的双手创造了我，形成了我}；而那位在\BibleBook{依撒意亚}{}中说：\BibleQuote{由母胎中形成我，为作祂仆人的上主这样说}\BibleRef{依 49:5}的，也表明了同样的道理。因此，就本性而言，他并不比我们有任何不同，尽管他在时间上先于我们，只要我们都由同一只手所造成和支持。那么，亚略派啊，如果你们对天主子也持这样的想法，认为祂也是如此存有和成为的，那么在你们看来，祂在本性上就与其他人毫无区别，只要祂过去也不存在，后来才成为，并且这名字是由于祂的德行而在创造时凭恩宠与祂结合的。因为按照你们的说法，祂自己就是其中之一，圣神在\BibleBook{}{圣咏}中论及这一类说：\BibleQuote{祂一发言，他们即成；祂一出命，他们即造。}若是这样，天主是通过谁发出命令来创造圣子的呢？因为必须有一个\OriginalTerm{圣言}{Word}，天主藉着祂发命，万物在祂内受造；但你们除了你们所否认的\OriginalTerm{圣言}{Word}之外，无法指出另一位，除非你们再捏造某种新的说法。\par

\Quote{是的，我们另有原因}；（这确实是我从前听\Person{欧瑟比}和他的同伙所采用的，）\Quote{我们之所以认为天主子比其他一切更具优越性，并被称为\OriginalTerm{独生子}{Only-begotten}，是因为唯有他分享父，而其他一切则分享子。} 因此，他们就这样费尽心机，像变换色彩一样改变和变化他们的说法；但这并不能使他们免受暴露，他们不过是属土的人，口出虚言，如陷入泥淖般沉溺于自己的幻想。\par

10. 因为如果祂被称为天主子，而我们被称为子之子，那么他们的虚构还算合理；但如果我们也被称为那位天主之子，而祂正是那位天主之子，那么我们也分享父，父说：\BibleQuote{我养育了儿女，并使他们长大}\BibleRef{依 1:2}。因为如果我们不分享祂，祂就不会说\Quote{我养育了}；但如果是祂自己生了我们，那么我们的父不是别的，正是祂。并且如前所述，子是否拥有更多且先受造，而我们则较少且后受造，这都无关紧要，只要我们都分享同一位父，并都被称为子。因为多与少并不表示本性不同，而是按照德行的实践加于各人；一人被置于十座城之上，另一人则五座；有些人坐在十二宝座上审判以色列十二支派；其他人则听到这样的话：\BibleQuote{我父所祝福的，你们来吧}，以及\BibleQuote{好！善良忠信的仆人}。然而，凭着这样的观念，难怪他们想象天主对于这样的子并非永远是父，这样的子也并非永远存在，而是从虚无中受造为受造物，并在其生成之前并不存在；因为这样的子不同于真天主子。\par

但坚持这样的教导并不符合虔敬，因为这更是撒杜塞人和\OriginalTerm{撒摩撒塔派}{Samosatene}的思想格调；那么，要说的就只剩下这一点了：天主子是按照另一种意义被称为子的，正如\Person{依撒格}是\Person{亚巴郎}之子；因为任何由某人生出，且非从外部加于他的，在本性上即是子，而这正是这名称所蕴含的意义。那么，圣子的生成是否出于人的情感呢？（或许他们的前辈们由于无知，也会急于提出这样的异议；）——绝非如此；因为天主不像人，人也不像天主。人是由物质受造的，且是可被情感左右的；但天主是非物质的，是无形的。如果圣经中对天主和人使用了同样的词语，但正如\Person{保禄}所嘱咐的，明眼人当研读圣经，并由此加以辨别，按照各个主题的本性来处置所写的内容，避免任何意义的混淆，既不以人的方式构想天主的事，也不将人的事归于天主。因为这无异于将酒掺水，并将凡火与神圣之火同置于\OriginalTerm{祭台}{altar}上。\par

11. 因为天主创造，而创造也被归于人；天主有\OriginalTerm{存在}{being}，而人也被称为\Quote{存在}，因他们从天主也领受了这恩赐。难道天主像人一样创造吗？或者祂的存在如同人的存在？绝无此意；我们是在一种意义上理解这些词语用于天主，在另一种意义上用于人。因为天主创造，是呼召那不存在的进入存在，无需任何条件；但人是加工已有的材料，先祈祷，从而从那位以祂\OriginalTerm{固有的圣言}{proper Word}创造万物的天主那里，获得制作的智慧。再者，人不能自存，被限定在空间内，且存立于天主的\OriginalTerm{圣言}{Word}中；但天主是自存的，包容万有，却不被任何所包容；按祂的善与能力，祂在万有之内，但按祂固有的本性，祂又在万有之外。因此，就像人创造不如天主创造，人的存在不如天主的存在一样，人的生育是一种方式，而圣子出于圣父是另一种方式。因为人的子女是父母的部分，由于身体的本质并非\OriginalTerm{单纯}{uncompounded}，而是处于流变状态，由部分组合而成；人在生育时失去其实体，又通过食物的补充重新获得实体。因此，人在时间中成为许多子女的父亲；但天主，由于没有部分，是圣子的父，既无\OriginalTerm{分割}{partition}也无\OriginalTerm{欲情}{passion}；因为\OriginalTerm{无质者}{Immaterial}既无\OriginalTerm{流溢}{effluence}，也无外来之灌注，不像在人中间那样；且因本性单纯，祂是独一圣子的父。这就是为什么祂是\OriginalTerm{独生子}{Only-begotten}，独自在父的怀里，也唯独被父宣认为出于祂，说：\BibleQuote{这是我的爱子，我所喜悦的}\BibleRef{玛 3:17}。祂也就是圣父的圣言，由此可以理解圣父\OriginalTerm{无欲情的}{impassible}且\OriginalTerm{不可分割的}{impartitive}本性，因为就连人的话都不是从情欲或分割受生的，何况天主的圣言呢？因此，祂也作为圣言，坐在圣父的右边；因为圣父在哪里，祂的圣言也在哪里；而我们，作为祂的工程，站在祂面前受审；祂受钦崇，因为祂是可钦崇的圣父之子，我们朝拜祂，宣认祂为主和天主，因为我们是受造物，异于祂。\par

12. 情况既是如此，就让他们中间愿意的人思考这事，好使某人可以用以下问题使他们羞愧：说天主的子嗣、属于祂本有的那位是从虚无来的，对吗？或者说，从天主而来的那位是附加上去的，这种想法合理吗，竟让人敢说圣子并非永恒存在？因为在此又显出圣子的\OriginalTerm{受生}{generation}超越并超乎人的思想，我们是在时间中成为自己儿女的父亲，因为我们自己先不存在，然后才成为存在；但天主，由于祂永远存在，永远是圣子的父。人类的起源是由类似的事使我们明白的；但既然\BibleQuote{除了父以外，没有人认识子；除了子和子所愿意启示的人以外，也没有人认识父}\BibleRef{玛 11:27}，因此，那些蒙子启示的圣作者们，就从可见之物给了我们某种形象，说：\BibleQuote{祂是天主光荣的辉耀，是祂本体的真像}\BibleRef{希 1:3}；又说：\Quote{因为在你那里有生命的泉源，借着你的光明我们才能看到光明}；当圣言申斥以色列时，祂说：\BibleQuote{你们离弃了智慧的泉源}\BibleRef{巴 3:12}；而这泉源就是那说：\Quote{他们离弃了我这活水的泉源}的。的确，这些比喻与我们想望的相比，是贫乏且非常暗昧的；但从中仍有可能理解一些超越人之本性的事，而不是把圣子的受生视为与我们的生育同列。因为有谁能想象光的\OriginalTerm{光辉}{radiance}曾不存在，以致敢说圣子并非永恒存在，或说圣子在受生前不存在呢？又有谁能把光辉与太阳分开，或设想泉源曾虚空无生命，以致狂妄地说：\Quote{圣子是从虚无来的}，这圣子却说：\BibleQuote{我是生命}\BibleRef{若 14:6}，或说\Quote{异于圣父的本质}，这圣子却说：\Quote{谁看见了我，就是看见了父？}因为圣作者们愿意我们如此理解，才给了这些比喻；当圣经含有这些形象时，从那些既非出自圣经、又无任何虔诚意义的其他说法中形成对我们主的观念，是不恰当且极其不虔敬的。\par

13. 因此，他们当告诉我们，他们是从哪位老师或是哪种传统中，得出关于救主的这些见解的？他们会说：\Quote{我们曾经读到，在\BibleBook{}{箴言}中：\BibleQuote{主造我于他化工之始}}；优西比乌及其同党常坚持这节经文，你也写信告诉我，如今这些人虽已被大量论证驳倒并挫败，却仍在各处散布这节经文，声称子是受造物之一，并将他归于受造之物。但在我看来，他们对此经文的理解也是错误的；因为它有着虔诚且极正统的意义，假如他们明白了，就不会亵渎荣耀的主了。若将上面所讲的与此经文相比较，他们就会发现二者之间差别甚大。凡是明理的人，谁不晓得，受造和制造之物是外在於制造者的？然而子，正如前述论证所表明的，并非外在，而是出自生他的父。因为人也建造房屋，也生儿子，没人会颠倒过来，说房屋或船是建造者所生的，而儿子是他所造所制作的；也不会说房屋是制造者的肖像，而儿子不像生他的人；他反而会承认儿子是父亲的肖像，房屋则是工艺作品，除非他心智混乱，失去理智。显然，圣经比任何事物都更了解万物的本性，它通过梅瑟论及受造物说：\BibleQuote{在起初天主创造了天地}\BibleRef{创 1:1}；而论及子时，它不是引入另一位，而是父自己说：\BibleQuote{我在晨星之前，从母胎中生了你}；又说：\BibleQuote{你是我的儿子，我今日生了你。}主也在\BibleBook{}{箴言}中论及自己说：\BibleQuote{早在山丘之前，他已生我}\BibleRef{箴 8:25}；对于受造和制造之物，若望说：\BibleQuote{万物都是借着他而造成的}\BibleRef{若 1:3}；但他宣讲主时却说：\BibleQuote{那在父怀里的独生子，已宣告了他。}那么，若是子，便不是受造物；若是受造物，便不是子；因为二者之间差别极大，子和受造物不能相同，除非他的本质被考虑为既出自天主，又外在於天主。\par

14. \Quote{那么，这节经文就没有意义了吗？}因为他们像一群蚊蚋般围著我们嗡嗡叫。绝非如此，它并非没有意义，而是有一个非常贴切的意义；因为说子也是受造的，这话确是真的，但这发生在他成为人的时候；因为受造属于人。任何人若不以研究神谕为次要之事，而是考察时间、特性与目的，并如此研读和思考他所阅读的，就会在神谕中发现这一意义被恰当地赋予。关于所说的时刻，他必会发现，虽然主常存在，但最终在时代的圆满中他成了人；他既是天主子，也成了人子。关于目的，他会明白，他愿意废除我们的死亡，便从童贞玛利亚那里取了一个身体；借著把这身体奉献给父，作为为众人的祭献，好能拯救我们众人，\BibleQuote{因为我们一生因怕死而受奴役}\BibleRef{希 2:15}。至于特性，这确实是救主的特性，但这是在取了一个身体并说\BibleQuote{主造我于他化工之始}\BibleRef{箴 8:22}时论及他的。因为正如永远存在并在父怀里，本是天主子的专有特性，同样，当他成为人时，\InnerQuote{主造了我}这话就适用于他。因为那时论及他，也说他饥饿、口渴、询问拉匝禄的坟墓、受苦、复活。正如当我们听到他说是主、天主、真光时，我们理解他是出自父；同样，听到\InnerQuote{主造了}、\InnerQuote{仆人}、\InnerQuote{他受苦}时，我们应将这不归于天主性，因为这不相关，而必须藉著他为我们所背负的肉体来解释：因为这些属于这肉体，而这肉体不是别人的，正是圣言的。如果我们想知道由此达成的目的，会发现如下：圣言成了血肉，为将这身体奉献为众人，而我们分享他的神，便得以\OriginalTerm{神化}{deified}，这是一份我们无法以其他方式获得的恩赐，除非他穿上了我们受造的身体，因为我们由此得名\Quote{天主的人}和\Quote{在基督内的人}。但正如我们领受圣神时，并不失去我们自身的本体，同样，主为我们成为人并背负身体时，仍旧是天主；因为他并未因身体的包裹而减损，反而使其神化并使其成为不朽的。\par

% structure: p0026 source=h2 target=section reason=relative-heading-rank; base=h2

\section{第4章——关于\InnerQuote{子}一词之公教含义的证明。\Emph{能力、圣言或理性，以及智慧，这些子的名称暗示著永恒；同样，父的\InnerQuote{泉源}称号亦如是。亚略派回答称，这些并不正式属于子的本质，而是赋予他的名称；天主有许多话语、能力等等。为何只有一个子和圣言等等。子的所有称号都在他内重合。}}

15. 至此，已足以揭露亚略异端的丑恶；因为，正如主所恩许，由他们自己的话，不虔就被归到他们身上。但如今，让我们主动进攻，要他们回答；因为现在是好时机，他们自己的立场失败了，轮到我们来质问他们；或许这能使悖逆者羞愧，并向他们指出他们从何处堕落了。我们从圣经得知，正如先前所说，天主子就是父的圣言和智慧。因为宗徒说：\BibleQuote{基督是天主的德能和天主的智慧}\BibleRef{格前 1:24}；\Person{若望}在说了\BibleQuote{圣言成了血肉}之后，立即补充说：\BibleQuote{我们见了祂的光荣，正如父独生者的光荣，满溢恩宠和真理}\BibleRef{若 1:14}，因此，圣言既是独生子，在这圣言和智慧内，天地及其中万物得以造成。我们由\Person{巴路克}得知，天主是智慧的泉源，因为以色列被指控离弃了智慧的泉源。那么，如果他们否认圣经，他们立刻就与基督徒的名号格格不入，可被众人恰当地称为无神论者和基督的仇敌，因为他们自己招来了这些名称。但如果他们与我们一致，承认圣经上的话是受天主默感的，就让他们敢于公开说出他们私下所思想的：天主曾一度无言无智慧；让他们的狂妄说：\Quote{曾有一个时候，祂不存在}，以及\Quote{在祂受生之前，基督不存在}；又让他们宣称，泉源并非从自身生出智慧，而是从外获得它，直至他们竟敢说：\Quote{圣子出自虚无}；由此将得出，那不再是泉源，而是一个水池，仿佛从外接水，却僭取了泉源之名。\par

16. 这充满了不虔，我认为但凡有一点理解力的人都不会怀疑。但既然他们咕哝着说，圣言和智慧不过是子的名称，我们就必须追问：如果这些仅仅是子的名称，那么子在它们之外必定另有实体。如果祂高于这些名称，则不能通过较低的名称来指称较高者；如果祂低于这些名称，那么祂必然在自身内具有此更尊贵称呼的原则；而这意味着祂的提升，这种说法与前面的不虔不相上下。因为那在父内，父也在祂内，曾说：\BibleQuote{我与父原是一体}\BibleRef{若 10:30}，人看见祂，就是看见了父；若说祂被任何外在的事物高举，真可谓是极端的疯狂。然而，当他们在这点上被驳倒，像\Person{欧瑟比}及其同党一样陷入极大困境时，他们还有最后的辩词，也就是\Person{亚略}在歌谣和他自己的\WorkTitle{塔利亚}中所编造的，作为一处新难题：\Quote{天主说了许多话；其中哪一个我们应称为子，为圣言，为父的独生子呢？}真是愚蠢，绝非基督徒！首先，使用这样的语言谈论天主，他们几乎将天主设想为人，一会儿说话，一会儿又用后来的话推翻先前的话，仿佛天主的一个圣言不足以按父的旨意形成万物，也不足以对万有施行天主眷顾。因为祂说许多话，正表明所有这些话的软弱，每一个都需要另一个的帮助。然而，天主只有一位圣言，这才是真道，既显示了天主的德能，也显示从祂而来的圣言之完美，以及如此相信者的虔诚领悟。\par

17. 但愿他们肯从自己这番言论出发宣认真理！因为如果他们一旦承认天主产生话语，他们就清楚地认识到祂是一位父；既承认这一点，就让他们想想：他们虽不愿将唯一的圣言归于天主，却想象祂是众多话语的父；他们虽不愿说天主根本没有圣言，却仍不宣认圣言就是天主的圣子——这是对真理的无知，也是对神圣圣经无经验。因为如果天主真是一位话语之父，为何受生者不是子呢？再者，除了祂的圣言，谁应是天主的圣子呢？因为话语不是众多，否则每一个都不完美；圣言却是唯一的，唯有祂是完全的，而且天主既唯一，祂的肖像也必须唯一，那就是圣子。因为天主的圣子，正如可从神圣的训示本身得知，祂自己就是天主的圣言、智慧、肖像、手和大能；因为天主的苗裔是唯一的，这些称号都是从父受生的标志。因为如果你说圣子，你就宣示了出自本性出于父的那位；如果你思念圣言，你就再次思念那出于祂且不可分离的；说到智慧，你又正是指那并非外来、而是出于祂并在祂内的；如果你称大能和手，你又是指那属于本质固有的；说到肖像，你便意指圣子；因为除了由祂所出的苗裔，还有什么能像天主呢？无疑，那藉着圣言而有的万有，都是\Quote{在智慧中立定}的，而那\Quote{在智慧中立定}的一切，都是藉着手造成的，并藉着圣子而有。对此我们有证据，并非来自外，而是出自圣经；因为天主亲自藉着先知\Person{依撒意亚}说：\BibleQuote{我的手奠定了大地，我的右手展开了诸天}\BibleRef{依 48:13}。又说：\BibleQuote{我要用我的手的荫庇遮蔽你，我正是用这手展开了诸天，奠定了大地的根基}\BibleRef{依 51:16}。\Person{达味}受此教导，知道主的手无非是智慧，便在圣咏中说，\BibleQuote{你以智慧造了万有；大地充满了你的受造物。} \Person{撒罗满}同样从天主领受了，说：\BibleQuote{主以智慧奠定了大地}\BibleRef{箴 3:19}，\Person{若望}知道那圣言就是手和智慧，便这样宣讲：\BibleQuote{在起初已有圣言，圣言与天主同在，圣言就是天主。圣言在起初就与天主同在。万有是藉着祂而造成的；凡受造的，没有一样不是由祂而造成的}\BibleRef{若 1:1}。那位宗徒看到手、智慧与圣言无非是圣子，就说：\BibleQuote{天主在古时曾多次并以多种方式，藉着先知对祖先说了话；但在这末期的日子里，祂藉着自己的圣子对我们说了话；天主已立祂为万有的承继者，并藉着祂创造了万世}\BibleRef{希 1:1}。又说：\BibleQuote{只有一个主，就是耶稣基督，万物藉着祂而有，我们也藉着祂而有}\BibleRef{格前 8:6}。他也知道那圣言、智慧、圣子本身就是父的肖像，便在\BibleBook{致哥罗森人}{书}中说：\BibleQuote{欣然感谢那使我们有资格在光明中分享圣徒福分的父，因为祂由黑暗的权势下救出了我们，并将我们移置在祂爱子的国内，我们且在祂内得到了救赎，即罪过的赦免。祂是不可见的天主的肖像，是一切受造物的首生者，因为在天上和在地上的一切，可见的与不可见的，或是上座者，或是宰制者，或是率领者，或是掌权者，都是在祂内受造的：一切都是藉着祂，并且是为了祂而受造的。祂在万有之先就有，万有都赖祂而存在。}因为正如万有是藉着圣言而受造，同样，因为祂是肖像，万有也是在祂内受造的。这样，凡将心思转向主的人，就不会绊倒在绊脚石上，反而会走向真理之光的灿烂中；因为这确实是真理的教义，尽管那些好争辩的人满怀怨恨地爆发，既不敬畏天主，面对自己的被驳倒也毫不羞愧。\par

% structure: p0030 source=h2 target=section reason=relative-heading-rank; base=h2

\section{第五章——维护大公会议的用语\Quote{\OriginalTerm{出自本质}{from the essence}}与\Quote{\OriginalTerm{同一本质}{one in essence}}。\Emph{反驳这些用语并非出自圣经；我们当注重意义胜于措辞；亚略派对圣经中\Quote{属于天主}一语的规避；他们规避一切解释，唯独接受大公会议所选定的，这些解释旨在驳斥亚略派的说法；抗议他们传递任何物质性的意义。}}

18. \Person{优西比乌}及其同伙先前已受到长时间详细审查，并且如我之前所说，他们自己定了自己的罪；于是他们签了字；在改变心意后，他们便安安静静地退休了；然而，如今这批人带着无宗教的狂妄新傲气，对真理昏头转向，极力指控大公会议，就让他们告诉我们，他们是从哪些圣经篇章学来的，或者是哪一位圣人教导了他们，使他们堆砌出这些短语：\Quote{\OriginalTerm{出于虚无}{out of nothing}}、\Quote{\OriginalTerm{在受生之前祂不存在}{He was not before His generation}}、\Quote{\OriginalTerm{曾有一时祂不存在}{once He was not}}、\Quote{\OriginalTerm{可变的}{alterable}}、\Quote{\OriginalTerm{先存}{pre-existence}}、\Quote{\OriginalTerm{凭意志}{at the will}}——这些都是他们用来嘲弄主的无稽之谈。因为真福\Person{保禄}在\BibleBook{致希伯来人}{书}中说：\BibleQuote{因着信德，我们知道普世是藉着天主的圣言所形成，以致看得见的是由看不见的所造成的}\BibleRef{希 11:3}。但圣言与普世毫无共通之处；因为祂是在普世之前就已存在的那一位，普世也是藉着祂而有的。并且在\WorkTitle{牧者}一书中也有记载（尽管这书不被列入\OriginalTerm{正典}{Canon}，但他们也拿这书来作证），其中写道：\Quote{首先当信天主是唯一的，祂创造、安排万物，并从虚无中引出万有}；但这仍然与圣子无关，因为这里论及的是那藉着祂而有的万有，祂自己与万有截然不同；因为我们不能将创造万有者也列入祂所造的万物之中，除非一个人疯狂到说建筑师就是他建造的房屋。\par

那么，为什么他们为了不虔敬的目的，自己发明了非圣经的用语，却指责那些在使用时怀着虔敬之心的人呢？因为不虔敬是绝对被禁止的，尽管有人试图用巧妙的措辞和似是而非的诡辩来掩饰它；但虔敬则被所有人承认是合法的，即使是用陌生的措辞提出，只要它们是以虔敬的态度使用的，并且希望它们成为虔敬思想的表达。正如前面所说，基督仇敌的那些卑劣用语过去和现在都充满着不虔敬；而大公会议针对它们的定义，如果仔细查考，就会发现它完全是真理的代表，尤其如果勤加注意产生这些用语的情由，这情由是合理的，其详情如下：——\par

19. 大公会议希望废除亚略派不虔敬的用语，而代之以圣经公认的言辞，即圣子不是从无中而是\Quote{从天主}，且是\Quote{圣言}和\Quote{智慧}，不是受造物或工程，而是父\OriginalTerm{本有的子嗣}{proper offspring}；\Person{欧瑟伯}及其同党，受他们根深蒂固的异端思想驱使，将\Quote{从天主？}这一短语理解为属于我们，好像就此而言天主的圣言与我们毫无差别，并且因为经上记载：\BibleQuote{只有一个天主，万有都出于他}（\BibleRef{格前 8:6}）；又说：\BibleQuote{旧事已过，都成了新的了；一切都是出于天主}（\BibleRef{格后 5:17}）。但教父们察觉到他们的诡计和他们不虔敬的狡猾，被迫更明确地表达\Quote{从天主}一语的含义。因此，他们写下了\Quote{\OriginalTerm{出自天主的本质}{from the essence of God}}，好使\Quote{从天主？}不被认为在圣子和受造物中是共同和同等的，而是使其他一切被承认为受造物，唯独圣言是出于父。因为虽然万有都说是出于天主，但这并不是在圣子出于他的那种意义上说的；因为至于受造物，\Quote{出于天主}之所以能说它们，是由于它们并非随意或自发地存在，也不是偶然产生——像那些将它们归于原子组合和类似结构元素的哲学家所说的那样——也不是像某些异端者所说的那样，有一位分开的造物主——也不是像另一些人所说，万有的构成出于某些天使；而是在于（天主既存在），是他借着他的圣言使万有从无中得以存在；但至于圣言，由于他不是受造物，唯有他既被称为且真是\Quote{出于父}。说圣子是\Quote{\OriginalTerm{出自父的本质}{from the essence of the Father}}正表显这层意义，因为这不附属于任何受造物。实际上，当保禄说\BibleQuote{万有都是出于天主}时，他立刻加上\BibleQuote{并只有一个主耶稣基督，万物借他而有}（\BibleRef{格前 8:6}），为要向众人显示，圣子异于这一切出于天主的事物（因为出于天主的事物，乃是借着他的圣子而成的）；而且他先前的话是指天主所造的宇宙说的，并非是说万有都像圣子那样出于父。因为其他事物既不像圣子，圣言也不与其他事物同类，因为他乃是万有的主与造物主；因此，神圣大公会议明确宣布他是\Quote{\OriginalTerm{出自父的本质}{from the essence of the Father}}，好使我们相信圣言有别于受造物的本性，唯独真实地出于天主；并且不给不虔敬之人留下任何口实。这就是大公会议写下\Quote{出自本质}的原因。\par

20. 再者，当主教们说\OriginalTerm{圣言}{Logos}必须被描述为圣父的\Emph{真能力}和\Emph{肖像}，在各方面与圣父完全相同、相像，且是不可改变的，是永远的，在祂内没有分离（因为圣言从没有不存在过，而是永远存在，与圣父永恒同在，如光的光芒），\Person{欧瑟伯}及其同党忍受了，不敢反驳，因为被驳得羞愧；然而他们被发现彼此私语、眨眼睛，说\Quote{相像}、\Quote{永远}、\Quote{能力}、\Quote{在祂内}，像以前一样，都是我们和子共有的，同意这些没有困难。至于\Quote{相像}，他们说经上指着我们写道：\BibleQuote{人是天主的肖像和光荣}\BibleRef{格前 11:7}；\Quote{永远}，则写道：\BibleQuote{因为我们这些活着的人，时常}\BibleRef{格后 4:11}；\Quote{在祂内}，则\BibleQuote{我们生活、行动、存在，都在祂内}\BibleRef{宗 17:28}；\Quote{不可改变}，则经上记载：\Quote{没有什么能叫我们与基督的爱隔绝}；至于\Quote{能力}，则蝗虫和蚱蜢被称为\Quote{能力}和\Quote{大军}\BibleRef{岳 2:25}，又说常论及民众，例如，\BibleQuote{主的军旅离开了埃及}\BibleRef{出 12:41}；还有别的，属天的，因为圣经说：\BibleQuote{万军的主与我们同在，雅各伯的天主是我们的保障}。的确，号称智者的\Person{阿斯特里乌斯}，已从他们那里学会并在著作中说了类似的话，在他之前的\Person{亚略}也已学会了，如前所述。但主教们看出了他们在这里的伪装，并且经上记载：\BibleQuote{图谋恶事的，心怀欺诈}\BibleRef{箴 12:20}，便再次不得不搜集圣经的意义，重新更清楚地重述并写下他们先前所说的话，即子与圣父\OriginalTerm{同性同体}{homoousion}：意思是要表明，子出于圣父，不仅相似，而是在相似上完全相同，并且显示子的相似和不可改变性，与归于我们的那种仿本不同，我们是因遵守诫命而修德得来的。因为彼此相似的身体可以分离并彼此远离，如同人类的子女与父母的关系（如论及\Person{亚当}和\Person{舍特}所记载的，他生了一个儿子，相似自己的模样\BibleRef{创 5:3}）；但既然子由圣父所生，不是依照人的性质，不仅相似，而且与圣父的本质不可分离，并且祂与圣父原是一体，正如祂自己所说，圣言常在圣父内，圣父也在圣言内，就像光芒与光的关系（因为这说法本身便指明的），因此大公会议理解这一点，恰当地写下了\OriginalTerm{同性同体}{homoousion}，好能击败异端者的乖僻，并显示圣言不同于受造之物。因为这样写后，他们立刻附加说：\Quote{但凡说天主子是从虚无中来的，或是受造的，或是可改变的，或是受造物，或是从别的本质来的，圣天主教会予以绝罚。}通过这样说，他们清楚表明，\OriginalTerm{出于本质}{ousia}和\OriginalTerm{同性同体}{homoousion}是摧毁那些不虔敬的口号的，如\Quote{受造的}、\Quote{受造物}、\Quote{受生的}、\Quote{可改变的}以及\Quote{祂在受生之前不存在}等。凡持守这些的，便是反对大公会议；凡不赞同\Person{亚略}的，必然持守并领受大公会议的决议，恰当地将其理解为指明光芒与光的关系，并由此获得真理的说明。\par

21. 所以，如果他们像其他人那样，以这些用语是陌生的为借口，请他们考虑大公会议如此写下的意义，并绝罚大公会议所绝罚的；然后，如果他们能，就让他们挑剔词句吧。但我深知，如果他们持守大公会议的意义，就会完全接受用来表达的那些词语；而如果他们想抱怨的是意义，那么所有人都明白，他们讨论措辞是徒劳的，不过是在为不虔敬寻找借口。这便是这些用语的缘由；但如果他们仍然抱怨这些用语不合圣经，这种抱怨本身就足以让他们被排斥，因为他们说空话，心志错乱。在这件事上，他们应当责备自己，因为他们开了先例，用不合圣经的言语向主宣战。然而，如果有人对此问题感兴趣，请让他知道，即使这些用语在圣经中并非字字句句都有，但如前所述，它们包含圣经的意义，并且表达出来，传达给那些听觉未损、能领悟虔敬教义的人。这一点你们要考虑，那些受教不良的人也应倾听。上文已经指出，且应相信为真，即\OriginalTerm{圣言}{Logos}出于圣父，是圣父独有的、本质的亲生之子。因为，除了从天主自己而来，我们还能设想子——那位智慧与圣言，万物藉之而受造——从何而来呢？然而，圣经也这样教导我们，因为圣父借\Person{达味}说：\BibleQuote{我的心灵涌溢美辞}，及\BibleQuote{在黎明之前，我由我胎中生了您}；子也向犹太人指明自己：\BibleQuote{如果天主是你们的圣父，你们必爱我，因为我是从天主出发而来的}\BibleRef{若 8:42}。又说：\Quote{除了那从天主来的那一位，没有人见过圣父；祂见过圣父。}再者，\Quote{我与圣父原是一体，}以及\Quote{我在圣父内，圣父在我内}等于说：\Quote{我出于圣父，且与圣父不可分离。}\Person{若望}在说：\BibleQuote{那在圣父怀里的独生子，他给我们详述了}时，讲的是他从救主那里学到的。此外，\Quote{在圣父怀里}还能暗示什么，不就是子由圣父的真实生育吗？\par

22. 因此，如果有人认为天主是复合的，如同\OriginalTerm{依附体}{accident}之在\OriginalTerm{本质}{essence}中，或有任何外在的包裹，并被环绕，或者好像有什么围绕祂而使其本质得以完备的东西，以至于当我们说\Quote{天主}或称呼\Quote{圣父}时，我们并非指那无形且不可理解的本质，而是指其周围的某个东西，那么他们尽管去抱怨大公会议宣布圣子出于天主的本质吧；但他们该想想，这样考虑他们是在说出两种亵渎的话；因为他们把天主变成了有形的，并且错误地说主不是真圣父的儿子，而是圣父周围某物的儿子。但如果天主是单纯的，也确实是单纯的，那么当我们说\Quote{天主}和称呼\Quote{圣父}时，我们并不是指称祂周围的什么，而是表示祂的本质本身。因为尽管理解天主的本质是什么是不可能的，然而只要我们理解到天主存在，并且如果圣经藉着这些称号指示祂，那么我们旨在指称祂而非任何其他，便称祂为天主、圣父和主。当祂说\BibleQuote{我是自有者}，以及\BibleQuote{我是上主天主}\BibleRef{出 3:14}时，或当圣经说\Quote{天主}时，我们由此所理解的，无非是暗示祂那不可理解的本质自身，以及那位被谈论者是\OriginalTerm{自有者}{He Is}。所以，听到天主圣子出于圣父的本质时，无人该感到惊愕；倒不如接受教父们的解释，他们以更明确但同义的语言，用\Quote{出于本质}取代了\Quote{从天主而来}。因为他们认为，说圣言\Quote{从天主而来}和说\Quote{出于天主的本质}是一样的，因为\Quote{天主}一词，如我已说过的，无非是指那位自有者的本质。那么，如果圣言不是像儿子出自父亲那样，是真正且由本性所生的，而是像受造物那样，因为是受造的，并且像\Quote{万物都由天主而来}那样，那么祂就既不是出于圣父的本质，圣子也不是根据本质而为子，而是由于德行，如同我们这些因恩宠而被称作子女的人。但如果祂唯独是由天主而来，如同真正的儿子，就如祂所是，那么就可以合理地说圣子出于天主的本质。\par

23. 再者，光和\OriginalTerm{光耀}{Radiance}的比喻具有这样的意义。因为圣人们并没有说，圣言与天主的关系有如由太阳的热所点着的火，那火通常又会熄灭，因为那是外在的工作，是造作者的一个受造物；但他们全都宣扬祂是\OriginalTerm{光耀}{Radiance}，以此表示祂出于本质，是固有的、不可分的，且与圣父合一。这也将确保祂的真正不变性和不可变动性；因为如果祂不是圣父本质固有的嗣子，这些又怎能属于祂呢？因为这必然也被用来印证祂与自己的圣父的同一性。我们的解释既如此虔诚，基督的仇敌也就不应该对\Quote{\OriginalTerm{同性同体}{One in essence}}感到惊愕了，因为这个术语也有其正确的意义和充足的理由。的确，如果我们说圣言出于天主的本质（因为依上述，这必须为他们所承认的说法），这除了意味着祂所由生的本质之真理和永恒，还能意味什么呢？因为这本质在种类上并无不同，免得它与天主的本质结合，如同某种外来的、不相似的东西。祂也不只是表面上相似，免得祂似乎在某个方面或全然在本质上是另一个，如同铜发光像金，银像锡。因为这些都是外来的、属于其他性质，在性质和能力上彼此分离，铜不是金所固有的，鸽子也不是从鸽子所生；虽然它们被认为相似，但在本质上却有差异。因此，如果圣子也是如此，就让祂像我们一样是受造物，而不是\OriginalTerm{同性同体}{One in essence}的吧；但如果圣子是圣言、智慧、圣父的肖像、光耀，那么按理祂必定是\OriginalTerm{同性同体}{One in essence}的。因为除非证明祂不是出于天主，而是一个在性质上不同、在本质上不同的工具，否则大公会议的教义就确实是健全的，其法令也是正确的。\par

24. 此外，关于此事，应摒弃所有物质性的联想；超越一切感官想象，让我们以纯洁的理智和纯粹的心灵，领悟圣子对圣父的真实关系，圣言对天主的专有关系，以及光辉对光的恒常肖似：因为正如\Quote{所生者}和\Quote{圣子}这两个词，所表达并旨在表达的不是任何属人的意义，而是合乎天主的，同样，当我们听到\OriginalTerm{同性同体}{one in essence}这一说法时，也不可陷入属人的感官，想象天主性具有分裂和划分，而要将思想导向非物质的事物，保持性体的一致和光的同一不可分割；因为这正是圣子对圣父的固有特性，并由此显明天主真是圣言的圣父。在此，光与其光辉的比喻也是恰当的。谁敢断言光辉与太阳不相像且为异物？相反，凡这样思想光辉与太阳的关系、光的同一性的人，谁不会满怀信心地说：\Quote{确实，光与光辉就是一，一者在另一者中显现，光辉在太阳内，以至于谁见此，即见彼？}然而，这样一种同一和本有的属性，那些信仰并正确识见的人应如何称呼，岂非同性同体的所生者？而天主的所生者，我们合宜地将其视为什么，岂非圣言、智慧、能力？若说这些与圣父相异，便是犯罪，甚至想象它们不与圣父永恒同在，也是恶行。因为借着这所生者，圣父创造了万物，并施其天主眷顾于万物；借着祂，圣父施行对人的爱，因此，正如所说，祂与圣父是一体；除非这些悖逆之人再作新论，说圣言的性体不是来自圣父而在祂内的光，仿佛圣子在内的光与圣父为一，但祂自身因是受造物，在性体上相异。然而这正是\Person{盖法}和\Person{撒摩撒特人}的信仰，教会已将其弃绝，但现在这些人却加以伪装；由此他们背离了真理，被宣称为异端者。因为如果祂完全分享来自圣父的光，为何祂不更是那其他存在所分享者，以免在祂与圣父之间引入中介？否则，便不再清楚万物是借圣子而受生，而是借那祂也分享者。若这就是圣言，圣父的智慧，在祂内圣父被启示和认识、创造世界，离了祂圣父一无所行，显然祂乃是从圣父而来者：因为所有受生之物都分享祂，如同分享圣神。祂既是如此，便不能从虚无而来，也不能是受造物，而是从圣父发出的本有的所生者，如同光辉来自光。\par

% structure: p0039 source=h2 target=section reason=relative-heading-rank; base=h2

\section{第六章 — 支持大公会议的权威证据。\Emph{\Person{德奥格诺斯图斯}；\Place{亚历山大}的\Person{狄奥尼修斯}；\Place{罗马}的\Person{狄奥尼修斯}；\Person{奥力振}。}}

25. 这就是在\Place{尼西亚}开会的人使用这些表述的含义。但下一步我们要证明，他们并非自己发明这些表述（因为这正是他们的一个借口），而是说他们从先辈领受的言语，这样也就堵住了他们的这个借口。因此，亚略派，基督的仇敌啊，你们要知道，博学者\Person{德奥格诺斯图斯}并不避讳\OriginalTerm{出于性体}{of the essence}这一说法，因为他在其\WorkTitle{纲要}的第二卷中，论及圣子时如此写道：——\par

\Quote{圣子的性体并非从外部获得的，也不是从虚无中产生的，而是从圣父的性体涌出，如同光的光辉，如同水的蒸汽；因为光辉或蒸汽，既不是水本身，也不是太阳本身，更不是异物；而是圣父性体的流出，然而，这流出并不遭受任何分割。因为正如太阳保持自身不变，并不因射出的光线而减损，同样，圣父的性体虽以圣子为其肖像，也不遭受改变。}\par

德奥格诺斯图斯经过先前的一种锻炼性研究之后，便以上述话语陈述了自己的观点。接下来，\Place{亚历山大}的主教\Person{狄奥尼修斯}，他曾撰文驳斥\Person{撒伯里乌}，并详细阐述了救世主因肉身而施行的救世工程，由此向撒伯里乌派证明，成为血肉的不是圣父，而是祂的圣言，正如\Person{若望}所说。但他因此被怀疑，说他认为圣子是一件受造且产生之物，并非与圣父\OriginalTerm{同性同体}{one in essence}；于是他致函与他同名的\Place{罗马}主教\Person{狄奥尼修斯}，为自己辩护，说这是对他的诬蔑。他向对方保证，自己从未称圣子为受造，反倒是承认圣子确实与圣父\OriginalTerm{同性同体}{one in essence}。他的言论如下：——\par

\Quote{我在另一封信中已驳斥了他们诬告我否认基督与天主\OriginalTerm{同性同体}{one in essence}的罪名。因为尽管我说过在\WorkTitle{圣经}中未找到这一术语，但我接下来的言论（他们并未注意）与这一信仰并不矛盾。例如，我举例人类生育显然是同质的，我观察到父母与其子女的区别无可否认地只在于非同一的个体，否则便无所谓父母与子女。至于那封信，如前所述，由于目前情况，我无法出示；否则我定会将原文，或更确切地说，整封信的抄本寄给你，若有机会，我仍会这样做。但我凭记忆确信，我曾引用过彼此同种类比：例如，从种子或根生出的植物，虽与所从出者不同，但在性体上却全然为一；又如，从泉源流出的溪流，获得了新的名称，因为泉源不称为溪流，溪流也不称为泉源，两者皆存在，而溪流是源自泉源的水。}\par

26. 天主圣言并非受造之物或受造物，而是圣父性体本有且不可分割的\OriginalTerm{所生者}{offspring}，正如大公会议所写，你们便可从\Place{罗马}主教\Person{狄奥尼修斯}的话中见到，他撰文反驳撒伯里乌派时，如此痛斥那些胆敢如此说的人：——\par

\Quote{接下来，我可以合理地转向那些分裂、割裂并毁灭天主的教会最神圣教义——\OriginalTerm{神圣君主制}{Divine Monarchy}的人，他们将其弄成三个权能和分割的\OriginalTerm{自立体}{subsistences}以及三个\OriginalTerm{神性}{Godheads}。我得知你们中间有些传道员和圣言的教师，在这教条上带头，他们可以说与\Person{撒伯里乌}的意见截然相反；因为他亵渎地说圣子就是圣父，圣父就是圣子，但他们则某种意义上宣讲三个天主，将神圣的\OriginalTerm{单一体}{Monad}分割成三个彼此相异且完全分离的自立体。因为宇宙的天主必然与\OriginalTerm{圣言}{Divine Word}结合，\OriginalTerm{圣神}{Holy Ghost}也必须寓居并住在天主内；因此，在一个首脑内，即宇宙的天主内，神圣的\OriginalTerm{圣三}{Triad}必须集合和联结。因为狂妄的\Person{马西昂}的教义就是要割裂并分割神圣君主制为三个本源——这是魔鬼的教导，不是基督的真门徒和热爱救主教训的人的教导。因为他们深知，\OriginalTerm{圣经}{divine Scripture}宣讲的是一位圣三，但\OriginalTerm{旧约}{Old Testament}和\OriginalTerm{新约}{New Testament}都没有宣讲三位天主。同样，必须谴责那些认为子是受造物，并认为主像那些真正受造之物一样形成的人；然而神圣的\OriginalTerm{谕言}{oracles}证实，子有适合祂的相称的\OriginalTerm{生发}{generation}，而非任何制作或创造。因此，说主无论如何是\OriginalTerm{手工艺品}{handiwork}，这不仅是普通的，甚至是最大的亵渎。因为如果祂成为圣子，就曾一度不是；但祂永远存在，即如果祂在父内，如祂自己所说，如果\OriginalTerm{基督}{Christ}是\OriginalTerm{圣言}{Word}、\OriginalTerm{智慧}{Wisdom}和\OriginalTerm{能力}{Power}（如你们所知，圣经这样说），而这些属性是天主的能力。那么，如果子曾形成，这些属性就曾一度不是；结果就会有一段时间天主没有这些属性；这是极其荒谬的。对于你们这些充满圣神、清楚知道说子是受造物所显出的荒谬的人来说，在这些问题上何必多言？依我看，这意见的作者们没有注意到这一点，在解释经文时，完全错过了真理，因为他们对神圣和先知圣经中那段经文的理解相悖，那句话说：\BibleQuote{主在祂行径之初，在造化万物之先，就造了我。}\BibleRef{箴 8:22}。因为\InnerQuote{祂创造}的含义，如你们所知，并非单一，我们必须将此处的\InnerQuote{祂创造}理解为\InnerQuote{祂立祂于祂所造的工程之上}，也就是\InnerQuote{圣子自己所造的工程}。这里的\InnerQuote{祂创造}绝不能理解为\InnerQuote{制造}，因为创造与制造不同。\BibleQuote{祂不是你的父，曾购买了你吗？祂不是制造了你，创造了你吗？}\BibleRef{申 32:6}——\Person{梅瑟}在\BibleBook{}{申命纪}的伟大诗歌中这样说。人们可以对那些轻率的人说：祂是受造物吗？那\InnerQuote{每样受造物的首生者，在曙光之前生于母胎}的，祂曾以智慧说：\BibleQuote{在丘陵还未存在之前，他就生了我。}\BibleRef{箴 8:25}。在神圣谕言的许多章节中，圣子被说是被生发的，却从未说是受造的；这明显定了那些误解主之生发的人的罪，他们竟敢称祂神圣而不可言喻的\OriginalTerm{生发}{generation}为制造。因此，我们既不能将美妙而神圣的\OriginalTerm{单一体}{Monad}分割为三个神性；也不能以\InnerQuote{受造物}之名贬低主的尊严和超越的荣威；而必须信仰全能的天主\OriginalTerm{圣父}{Father}，和祂的子耶稣基督，以及圣神，并持守圣言与宇宙的天主结合。因为祂说：\BibleQuote{我与父是一体}；又说：\BibleQuote{我在父内，父在我内}。因为这样，神圣的圣三和关于神圣君主制的宣讲必得以保全。}\par

27. 关于\OriginalTerm{圣言}{Word}与\OriginalTerm{父}{Father}的永恒共存，以及祂并非来自另一\OriginalTerm{本质}{essence}或\OriginalTerm{自立体}{subsistence}，而是父所专有的，正如大公会议中的主教们所说，你们还可以从勤勉的\Person{奥利振}那里听到。凡他作为探讨和练习所写的内容，任何人都不要当作表达他自己的意见，而是代表在研究中争辩的各方，但他明确宣告的，正是这位勤勉者自己的意见。在他对\OriginalTerm{异端者}{heretics}作了（可以说是）开场白之后，他立刻介绍自己的信仰，如下：——\par

\Quote{如果存在\OriginalTerm{不可见天主的肖像}{Image of the Invisible God}，那就是不可见的肖像；我且大胆加上一句，既是\OriginalTerm{父}{Father}的相似者，就从来不曾没有。因为那位按照\Person{若望}被称为\OriginalTerm{光}{Light}的天主（\BibleQuote{天主是光}），何时曾没有自己本有荣耀的\OriginalTerm{光辉}{radiance}，以致人可以竟敢宣称\OriginalTerm{子}{Son}的存在有起始，好像祂从前不存在？但父那不可言说、无名且无法表达的\OriginalTerm{自立体}{subsistence}，那\OriginalTerm{表达和圣言}{Expression and Word}，以及认识父的那位，何时曾不存在？因为敢于说\InnerQuote{曾有一时子不存在}的人，他必须明白自己是在说\InnerQuote{曾有一时智慧不存在}、\InnerQuote{圣言不存在}、\InnerQuote{生命不存在}。}\par

他在别处又说道：——\par

\Quote{但如果因为我们悟性的软弱，我们竟在力所能及的范围内剥夺天主那永远与祂共存的\OriginalTerm{独生圣言}{Only-begotten Word}，以及祂所喜悦的\OriginalTerm{智慧}{Wisdom}，那便不是无罪的，也不是没有危险的；否则，就必须设想天主并非永远拥有喜乐。}\par

看，我们正证明这观点是从\OriginalTerm{教父}{father}传到教父的；但你们，现代的\OriginalTerm{犹太人}{Jews}和\Person{盖法}的门徒啊，你们能举出多少教父来支持你们的说法？没有一个有悟性和智慧的人；因为所有人都厌恶你们，唯有魔鬼不厌；在这\OriginalTerm{背教}{apostasy}中，只有他是你们的父，他既在起初将这不信的种子撒在你们里面，现在又劝诱你们诽谤\OriginalTerm{大公会议}{Ecumenical Council}，因为它所记录的并非你们的教义，而是从起初那些亲眼看见并作\OriginalTerm{圣言}{Word}仆役的人所传给我们的。因为大公会议所书面宣认的信仰，正是\OriginalTerm{天主教会}{Catholic Church}的信仰；为确认这点，\OriginalTerm{蒙福的教父}{blessed Fathers}们在谴责\OriginalTerm{亚略异端}{Arian heresy}时如此表达了；而这也是这些人转而诽谤大公会议的主要原因。因为困扰他们的并非那些措辞，而是那些措辞证明了他们是\OriginalTerm{异端者}{heretics}，并且狂妄超过其他异端。\par

% structure: p0051 source=h2 target=section reason=relative-heading-rank; base=h2

\section{论亚略派的标记\Quote{非受生}。\Emph{此词后来为他们所采用；其原因；该词的三种含义。第四种含义。非受生指天主与其受造物相对，而非与其子相对；圣经所用之称号为父；结论。}}

28. 这正是在那时，当他们语词之虚谬本质已被揭露，他们自此受到不敬虔之指控后，他们之所以转去从希腊人那里借来\OriginalTerm{非受生}{Unoriginate}一词的原因；在它的遮掩下，他们好将那天主圣言——万物原是藉祂而造成的——算入受造之物中；他们在不敬神中如此不知羞耻，在对主的亵渎中如此顽梗。如果这种无耻源于对该词的无知，他们本应向那些将此词传授给他们的人学习，那些人毫不迟疑地说：即使是理智——他们从至善推演而出——以及从理智而出的灵魂，虽然它们各自的起源为人所知，却仍然是非受生的，因为他们明白，如此说并不贬损其他事物由之而来的那第一起源。既如此，他们就该自己也这样说，否则便绝不要谈论他们不知之事。然而，若他们认为自己熟悉这题目，那么就必须质问他们；因为这表达不出自圣经，而他们是在争论，如在别处一样，为那些无圣经根据的立场。正如我已说明了那理由与意义，即大公会议和先前的教父们依据圣经论及救主之处，界定并公布\Quote{属于本质}和\Quote{同一本质}的理由和意义；那么现在让他们，若他们能，就自己方面回答，是什么引导他们使用这无圣经依据的语词，以及他们在何种意义上称天主为非受生的？事实上，我听说，这名称有不同意义；哲学家们说，它首先指\Quote{尚未生成，但可能生成之物}；其次指\Quote{既不存在，也不能生成之物}；第三指\Quote{确实存在，但既非受生，也无存在之起源，而是永恒且不可毁灭的}。现在，或许他们会愿意略过前两种意义，因为随之而来的荒谬；因为按照第一种，那已经生成之物，和那期待生成之物，都是非受生的；而第二种更为荒谬；因此他们会采用第三种意义，并以此意义使用该词；尽管在这里，即使按此意义，他们的不敬依然毫不逊色。因为，若他们以非受生意指那无存在起源、既非受生亦非受造、而是永恒之物，且说天主圣言与此相反，谁不明白这些天主仇敌的诡计？谁不愿意用石头砸死这样的疯子？因为，当他们羞于再次提出那些他们杜撰且已被谴责的前述说法时，这些恶人便通过他们所谓的非受生之说，另辟蹊径以表达同样的意思。因为，若圣子属于受造之物，那么，祂也是从虚无而成为的；若祂有存在之起源，那么，祂在受生之前就不存在；若祂不是永恒的，就曾有一时祂不存在。\par

29. 若这便是他们的想法，他们就该以自己的说法表明其异端，而不该将他们的乖谬藏在\OriginalTerm{非受生}{Unoriginate}的伪装之下。然而相反，这心怀恶念的人凡事都以诡诈而行，如同他们的父魔鬼；因为正如魔鬼试图假借他人形貌进行欺骗，这些人便推出\Quote{非受生}一词，以便他们可以佯装恭敬地谈论天主，却在暗中怀藏对主的亵渎，并在遮掩之下将它教导他人。然而，当这诡辩被识破，他们还有什么可剩下的呢？\Quote{我们找到了另一个，}作恶者说；于是进一步对他们已说过的加以补充，就是\Quote{非受生}意指那没有存在之作者，而它本身对于受造之物来说便是如此关系。忘恩负义，且实在地对圣经充耳不闻！他们做一切事、说一切话，不是为了荣耀天主，而是为了羞辱圣子，竟不知羞辱圣子的，就是羞辱圣父。因为首先，即使他们以这种方式称呼天主，圣言仍未证明属于受造之物。因为反过来，作为圣父本质的子嗣，祂自永恒便与圣父同在。因为这子嗣之名并不减损圣言的本性，\Quote{非受生}也不在与圣子的对比中获得其意义，而是与藉着圣子而成的事物相对；正如有人称呼一位建筑师，且称他为房屋或城市的建造者时，并不能因这个称呼而暗指他亲生的儿子，而是因他在工作中所展示的技艺和知识而称他为工匠，以此来指明他并非如同他所造的事物，并且他既知道建筑者的本性，也知道他所生养的儿子有别于他的作品；在面对他儿子时称他为父，而在面对他的作品时，称他为创造者和制造者；同样，那在此意义上说天主是非受生的人，是从祂的工程来称呼祂，不仅指明祂并非受生，且指明祂是那些受生之物的制造者；但他也明白，圣言与受生之物不同，且是圣父的独生子，万物都是藉着祂而成就并立定的。\par

30. 同样，当先知们称天主为\OriginalTerm{统御万有者}{All-ruling}时，他们并非这样称呼祂，仿佛圣言被包含在那万有之中（因为他们知道子不同于\OriginalTerm{受造之物}{originated}，且按祂与父的相似，祂自己就是万有的主宰）；而是因为祂是通过子所创造的万物的统治者，并将万有的权柄交给了子，及至交了之后，祂自己又藉着圣言再次成为万有的主。再者，当他们称天主为\OriginalTerm{万军之主}{Lord of the powers}时，他们说这话并非仿佛圣言是那些大能者之一，而是因为祂既是子的父，也是藉着子而生成的那些大能者的主。因为同样，圣言既在父内，也就是万有的主和主宰；因为凡父所有的一切，都是子的。这就是这类称号的含义，同样，人若愿意，也可称天主为\OriginalTerm{非受生者}{unoriginated}；但这并非表示圣言属于\OriginalTerm{受造之物}{originated things}，而是因为，如我先前所说，天主不仅不是受生的，而且是通过祂自己的圣言成为那些\OriginalTerm{受造之物}{things which are so}的创造者。因为父虽被如此称呼，圣言仍是父的\OriginalTerm{肖像}{Image}，与祂\OriginalTerm{同性同体}{one in essence}；既是祂的肖像，就必与\OriginalTerm{受造之物}{things originated}及一切有别；因为祂是谁的肖像，就具有谁的特性和模样：因此，凡称父为非受生者和全能者的人，在那非受生者和全能者中领悟到祂的圣言和智慧，也就是子。但这些可惊而急于不信的人，抓住了非受生者这一词语，并非出于对天主的尊敬，而是出于对救主的恶意；因为他们若真注重尊敬和虔敬的语言，那么承认并称呼天主为父，要比给祂这个名称更正确且更好；因为称呼天主为非受生者，如我先前所说，他们是从\OriginalTerm{受造之物}{things which came to be}来称呼祂，且仅视之为创造者，好使他们能按自己的喜好暗示圣言是工程；但称呼天主为父的人，在其中同时意指祂的子，并且不可能不知道，既有子，万物就是借着这个子而受造的。\par

31. 因此，从子来指称天主并称祂为父，比仅凭祂的工程来命名祂、称祂为非受生者要准确得多；因为后一个词指涉的是那些藉由圣言、按天主的旨意而生成的工程，而父的名称则指出来自祂本质的本有子嗣。而且，圣言超越受造之物有多少，称天主为父就比称祂为非受生者超越更多；因为后者既非出自圣经，又令人疑虑，且歧义甚多；而前者简单明了、合乎圣经、更为准确，且独自暗指子。\InnerQuote{非受生者}是那些不认识子的希腊人的用语；但\InnerQuote{父}却是我们的主所承认并赐予的；因为祂知道自己是谁的子，曾说：\BibleQuote{我在父内，父在我内}\BibleRef{若 14:9} 又说：\BibleQuote{看见我的，就是看见了父}；\BibleQuote{我与父原是一体}；但从未见祂称父为非受生者。此外，当祂教导我们祈祷时，祂不说：\InnerQuote{你们祈祷时，要说：非受生的天主啊}，而是说：\InnerQuote{你们应当这样祈祷：\BibleQuote{我们在天的父！}\BibleRef{玛 6:9}}。祂也愿意我们信仰的纲要具有相同的取向。因为祂命令我们受洗，不是归于非受生者与受造者的名，也不是归于\OriginalTerm{非受造者}{Uncreate}与\OriginalTerm{受造物}{Creature}的名，而是归于父、子、圣神的名，因为藉着这样的入门圣事，我们也真正成了儿子，并使用父的名称，从这一名称认出在父内的圣言。但若祂愿意我们称祂自己的父为我们的父，我们不可因此按本性衡量自己与子同等，因为父是因着子才被我们如此称呼的；因为圣言承担了我们的肉体，并在我们内生成，所以因着在我们内的圣言，天主才被称为我们的父。因为我们内的圣言之神，通过我们称呼祂自己的父为我们的父，这正是宗徒的意思，当他说：\BibleQuote{天主派遣了祂子的神，进入我们心中，呼喊：阿爸，父啊！}\BibleRef{迦 4:6}。\par

32. 但或许在\Quote{非受生者}这个术语上被驳倒后，他们会按其邪恶的本性说：\Quote{关于我们的主、救主耶稣基督，也理应从圣经中陈述那里所载关于祂的事，而不应引进圣经之外的词语。}是的，我也说理应如此；因为从圣经得来的真理标记，比从其他来源得来的更为精确；但\Person{欧瑟伯}及其同党恶劣的性情和反复无常且狡诈的不虔，迫使主教们，正如我先前所说，更明确地发表了那些推倒其不虔的词语；而公会议所写下的，已显明具有正统的意义，同时\OriginalTerm{亚略派}{Arians}则在言语上败坏、性情上邪恶已得到证明。\OriginalTerm{非受生者}{Unoriginate}这个术语，自有其意义，且容人虔诚地使用，他们却按自己的想法和意愿，用来侮辱救主，这一切都是为了像巨人一样，执拗地维持他们与天主的战斗。但是，当他们提出这些先前的话语时既未能逃脱谴责，同样，当他们错误地理解\OriginalTerm{非受生者}{Unoriginated}——它本身容许被妥善而虔诚地使用——时，他们便被揭穿了，在众人面前蒙羞，其异端到处被禁绝。以上，我已经尽我所能，述说了以往在公会议中所发生的事，以作解释；但我知道基督的仇敌中那些争强好胜的人，即使在听了这些之后也不会改变，而是会不断寻找别的借口，而在这之后又再寻找别的借口。因为正如先知所说：\BibleQuote{雇士人岂能改变他的肤色？豹岂能改变它的斑点？}\BibleRef{耶 13:23}，那时，那些受教于不虔的人才会愿意虔诚地思想。然而，亲爱的，你收到这封信后，请自己阅读；若蒙嘉纳，也读给恰好与你在一起的弟兄们听，好叫他们听见了，也能欢迎公会议追求真理的热忱及其意义之精确；并谴责基督的仇敌——\OriginalTerm{亚略派}{Arians}，以及他们为维护其不虔的异端而在自己中间苦心营造的那些徒劳的借口；因为光荣、崇敬和敬拜，当归于天主圣父，连同祂同存之圣子圣言，以及全圣而赋予生命的圣神，从现在直到永世无穷之世。阿们。\par

\SourceRecord{Translated by John Henry Newman.；Nicene and Post-Nicene Fathers, Second Series；Vol. 4.；Edited by Philip Schaff and Henry Wace.；Buffalo, NY: Christian Literature Publishing Co.,；1892.；<http://www.newadvent.org/fathers/2809.htm>.}
